\section{Related work} 

\textbf{Generalization of weight decay.} 
In classical statistics and machine learning, improving model generalization via weight decay has long been an important research topic \cite{doi:10.1137/0109031,Hoerl,a92f3c16-7c6e-31d3-b403-82d2b0a469e4}. \citet{10.5555/2986916.2987033} show in linear networks that weight decay suppresses irrelevant components of the weight vector and mitigates the effect of noise in the target values, thereby significantly improving generalization performance. More generally, \citet{ShalevShwartz2014UnderstandingML} and \citet{Zhang2016UnderstandingDL} provide theoretical generalization analyses supporting the regularization effect of weight decay, although the latter noted that its improvement on test accuracy can be limited in practice. Building on adaptive optimization, \citet{Loshchilov2017DecoupledWD} highlight the fundamental difference between weight decay and $L_2$-regularization in Adam, demonstrating that AdamW achieves more effective regularization while simplifying hyperparameter tuning. From the Bayesian posterior aspect, \citet{Zhou2024TowardsUC} prove the convergence of AdamW and justify its generalization advantages over Adam and $L_2$-Adam.

\textbf{Parameter settings related to weight decay.} 
\citet{Loshchilov2017DecoupledWD} explicitly highlights the coupling between the learning rate $\eta$ and the weight decay coefficient $\lambda$ in their study of Adam. For another parameter, \citet{blake2025umupunitscaledmaximalupdate} observe that, under fixed batch size and weight decay, the optimal learning rate decreases as the dataset size $N$ increases. This finding is broadly consistent with the results of \citet{wang2025setadamwsweightdecay}, who shows that larger datasets require a longer effective time scale $T$, which can be achieved by either reducing the weight decay or decreasing the learning rate. \citet{bergsma2025powerlinesscalinglaws} further demonstrates that the optimal weight decay coefficient scales linearly with the batch size $B$. Compared to directly tuning the learning rate, adjusting $\lambda$ provides improved stability in practice. In addition, \citet{fan2025robustlayerwisescalingrules} and \citet{wang2025setadamwsweightdecay} have shown that weight decay is critical to transfer efficacy. Similar to the conclusions given by \citet{wang2025setadamwsweightdecay}, \citet{dangelo2024needweightdecaymodern} observe that under under-training conditions, the optimal generalization performance of SGD often approximates $\eta \propto \frac{1}{\lambda}$. 

The work most similar to ours is \cite{wang2025setadamwsweightdecay}. They investigate the parameter settings of the AdamW algorithm by one-to-one mapping from the EMA timescale to the weight decay hyperparameter. However, in this work, we aim to study the coupling effects between $\lambda$ and other optimization parameters by deriving explicit expressions of weight decay under generalization-driven analysis.

